\section{Theoretical Framework}

\subsection{Modern IoT Backends}

\subsubsection{Backend-as-a-Service (BaaS): The Role of Supabase in IoT}
A Backend-as-a-Service (BaaS) is a cloud-based platform that provides ready-made backend components such as databases, authentication, file storage, and APIs so that developers can focus mainly on the application logic and user interface instead of server management. In practice, this means the platform hides (abstracts) the complexity of provisioning servers, configuring databases, and maintaining security updates. Supabase is an example of a BaaS that is built directly on PostgreSQL, meaning that all data is stored in a relational Postgres database while the platform adds services like authentication, real-time subscriptions, and serverless edge functions around it.\cite{blog_elest_supabase_self_host, supabase}

Supabase exposes a Postgres-native backend, where the database schema defines the structure of the API, and auto-generated REST and real-time APIs are created on top of the same relational data model. Real-time subscriptions are implemented using WebSockets, which are long-lived network connections that allow a server and client to send messages to each other at any time without opening a new HTTP request for each update. When a row in a watched table changes, Supabase’s Realtime service streams an event over the WebSocket so that clients can immediately react, for example by updating a dashboard or triggering an alert.\cite{bejamas_serverless_supabase, supabase, blog_elest_supabase_self_host}

The serverless paradigm means that developers deploy small units of code (functions) without managing the underlying servers; the BaaS platform automatically scales these functions up or down based on incoming requests. In this model, billing and capacity planning are usually based on usage (invocations, execution time, or bandwidth) rather than fixed server instances, which is attractive for IoT workloads with fluctuating traffic. A unified backend refers to the idea that authentication, database access, real-time messaging, and serverless business logic are all offered by a single platform with consistent APIs and permissions, instead of being assembled from multiple loosely integrated services. This unified view simplifies backend development for IoT applications that must combine sensor ingestion, user management, and real-time visualization.\cite{deepwiki_supabase_core, supabase, blog_elest_supabase_self_host, bejamas_serverless_supabase}

\textbf{[INSERT DIAGRAM: Supabase Realtime Architecture Flow]}

The Supabase realtime architecture can be viewed as a flow where IoT clients or backend services write data into PostgreSQL, the Realtime service listens to database changes, and interested clients maintain WebSocket subscriptions to be notified whenever relevant rows are inserted, updated, or deleted.\cite{blog_elest_supabase_self_host, bejamas_serverless_supabase}

\subsubsection{Time-Series Databases: Why standard SQL fails for sensor streams}
Traditional relational databases are optimized for transactional workloads where each query touches a relatively small set of rows, such as updating a single user record, rather than continuously ingesting and querying millions of timestamped sensor readings per day. High-velocity ingestion describes scenarios where large amounts of data arrive in short intervals, for example multiple measurements per second from many devices; under such conditions, a conventional relational engine can become a bottleneck because each insert must update multiple indexes and write to disk frequently.\cite{tigerdata_tsdb_vs_rdbms, cratedb_tsdb_vs_rdbms}

Most relational systems rely on B-tree indexes, which are tree structures that help locate rows quickly by sorting keys and storing them in multiple linked pages. When the index and data no longer fit into memory, random inserts into a large B-tree force the database to load and write many different pages, increasing disk I/O and slowing down both inserts and queries. For sensor workloads, data forms a continuous stream, meaning new records keep arriving over time with timestamps that grow (mostly) monotonically, and queries often scan over ranges of time rather than single rows. As the volume grows, query degradation occurs: time-range queries become slower because they must process more rows and more index pages, and maintenance operations such as vacuuming and index rebalancing also add overhead. Retention issues appear when old sensor data is kept indefinitely, leading to storage bloat and longer query times; without dedicated retention mechanisms, deleting large ranges of old data can be expensive and disruptive.\cite{cratedb_tsdb_vs_rdbms, tigerdata_tsdb_vs_rdbms}

\subsubsection{The TimescaleDB Extension: Hypertables and Chunking explained}
TimescaleDB is an extension to PostgreSQL that turns it into a time-series database by adding new concepts on top of the existing relational engine while keeping full SQL compatibility. An extension in this context means a packaged module that plugs into PostgreSQL to add new functionality (such as new data types or storage methods) without replacing the core database server. A hypertable is a logical table abstraction offered by TimescaleDB that automatically partitions time-series data into many smaller physical tables, called chunks, based on time (and optionally space, such as device ID), while allowing the user to query it as if it were a single regular table.\cite{tigerdata_caggs_retention, so_timescaledb_size_retention, timescaledb_retention_caggs}

A chunk is one of these underlying physical partitions that stores a subset of the hypertable’s data, typically corresponding to a time interval (for example, one day or one week of readings). TimescaleDB uses automatic chunking to exploit temporal locality, which means that most queries and inserts operate on recent time intervals and can therefore hit a small number of chunks, reducing index and I/O overhead compared to a single huge table. This design enables high-performance inserts because new data tends to go into the most recent chunk, which remains hot in memory and can be written sequentially. Compression features allow older chunks to be compressed to save storage while still being queryable, which is useful for long-term archives of sensor data.\cite{so_timescaledb_size_retention, tigerdata_tsdb_vs_rdbms, tigerdata_caggs_retention}

Continuous aggregates are materialized views that TimescaleDB keeps updated over time by periodically summarizing raw data into aggregated form (for example, hourly averages), which significantly speeds up analytical queries over long periods. A retention policy in TimescaleDB defines rules for automatically dropping old chunks from hypertables or continuous aggregates once they are older than a configured time interval, allowing the system to control storage growth without manual cleanup. Because TimescaleDB is built as an extension, it preserves SQL compatibility: existing PostgreSQL tools, queries, and drivers work largely unchanged, while the time-series optimizations remain mostly transparent to the application developer.\cite{dbt_timescaledb_retention, timescaledb_retention_caggs, tigerdata_docs_caggs_retention, tigerdata_caggs_retention}

\textbf{[INSERT DIAGRAM: Hypertable vs. Normal Table Storage]}

The hypertable versus normal table diagram conceptually contrasts a single, monolithic table with one where data is internally split into multiple time-based chunks, illustrating how queries over recent data can target only a few chunks instead of scanning the entire dataset.\cite{tigerdata_tsdb_vs_rdbms, tigerdata_caggs_retention}

\subsubsection{Object Storage and Media Handling in Microservices}
\begin{itemize}
    \item Unstructured data storage, Supabase Storage, S3 compatibility, scalable object storage, media asset management, separation of concerns (DB vs. Storage), CDN caching strategies, secure file access policies, handling large binary files (blobs)
    \item Source: \url{https://github.com/supabase/storage}
\end{itemize}

\subsection{Microservices in IoT}

\subsubsection{The Gateway Pattern: Bridging MQTT and the Database}
A microservice is a small, independently deployable service that implements a focused piece of functionality and communicates with other services over the network. In a microservice architecture, many such services collaborate to form the full system, and each can be scaled, updated, or restarted without affecting the others; this is known as service isolation, because failures or performance issues in one service are less likely to propagate to others. Fault isolation is the related property that when a service crashes or misbehaves, the impact on the overall system is limited, for example by using retry logic or circuit breakers at the boundaries.\cite{deepwiki_supabase_core}

The gateway pattern is a design pattern where a dedicated service acts as an entry point between one protocol domain and another, such as between MQTT messages from IoT devices and database writes in the backend. In an IoT context, a gateway microservice subscribes to MQTT topics, validates and transforms the incoming payloads, and then writes them into the time-series database, possibly enriching them with metadata or applying basic business rules. By centralizing protocol translation and validation, the gateway pattern improves scalability, because multiple gateway instances can run in parallel, and supports fault isolation, because failures in MQTT handling do not directly affect downstream services that query the database.\cite{microservices_io_gateway}

\subsubsection{Service Inter-communication: Synchronous vs. Asynchronous Patterns}
Synchronous communication means that one service sends a request to another and waits for an immediate response before continuing, as is typical for HTTP request–response interactions. This pattern is simple to reason about but introduces tight coupling and latency sensitivity, because the caller is directly dependent on the callee’s availability and response time. Asynchronous communication means that a service sends a message or event without waiting for an immediate reply, often by using message queues, publish–subscribe systems, or event streams; the sender can continue its work while the receiver processes the message later.\cite{gfg_microservices_communication}

An event-driven architecture is a style in which services react to events—discrete messages that describe something that happened, such as “sensor reading stored” or “threshold exceeded”—rather than continuously polling for changes. Coupling describes how strongly services depend on each other’s internal behavior and timing; synchronous RPC-style communication tends to create higher coupling, while asynchronous, event-driven patterns reduce coupling by decoupling senders and receivers in time and often in knowledge of each other’s identities. In IoT backends, choosing between synchronous and asynchronous inter-service communication influences latency (how quickly results are visible), reliability (how well the system tolerates temporary failures), and scalability (how easily throughput can be increased by adding more instances).\cite{gfg_microservices_communication}

\subsection{Edge Computing Patterns}

\subsubsection{Cloud, Fog, and Edge Computing}
Cloud computing refers to running applications and storing data in large, centralized data centers that are accessed over the internet, usually operated by hyperscale providers. Edge computing, in contrast, moves computation and decision-making closer to where data is produced, such as on gateways, local servers, or even directly on devices, to reduce network usage and response time. Fog computing is an intermediate layer between the edge and the cloud, typically consisting of regional or on-premises infrastructure that aggregates data from many edge nodes and performs preprocessing before forwarding summaries to the cloud.

Latency is the time delay between an input (for example, a sensor reading) and the observable reaction of the system (for example, an actuator change or dashboard update); moving computation closer to the data source generally reduces latency because fewer network hops and less bandwidth are required. Resilience describes the ability of a system to continue operating acceptably in the presence of failures or network partitions, which is often improved when some processing is performed at the edge or fog layer so that local behavior can continue even if the cloud is temporarily unreachable. Processing location trade-offs in IoT therefore involve balancing low latency and high resilience at the edge against the centralized scalability and advanced analytics capabilities of the cloud.

\subsubsection{Thick Edge vs. Thin Edge}
A thick edge architecture places significant logic and decision-making directly on the edge device or gateway, such as performing local anomaly detection, threshold checks, and actuation decisions without always consulting the cloud. In this model, the device may store recent history, apply filtering or aggregation, and only send summarized or exceptional data upstream, which can reduce bandwidth usage and allow real-time reactions even when connectivity is limited. A thin edge architecture, in contrast, keeps the edge device relatively simple, mainly responsible for collecting raw sensor values and forwarding them reliably to the backend, where most processing, analytics, and decision logic are executed.

The trade-offs between thick and thin edge designs involve complexity, maintainability, and reliability: thick edge devices are more capable but harder to update and test at scale, while thin edge devices are easier to manage but rely more heavily on network connectivity and backend availability for timely decisions. For real-time plant monitoring, the choice determines whether actions such as irrigation adjustments happen locally based on on-device rules or centrally in the cloud based on aggregated analytics and historical trends.

Modern IoT architectures therefore combine BaaS-style backend services, time-series optimizations such as hypertables and retention policies, microservice-based decomposition with gateway patterns, and carefully chosen processing locations across edge, fog, and cloud layers. These concepts enable systems that can ingest high-velocity sensor streams, store and query them efficiently, and remain scalable and resilient under changing workloads. In the next section (3.3), these theoretical principles are applied to the concrete PlantUp architecture, showing how the abstract patterns presented here are realized in a practical smart gardening system.
